
\documentclass{article}
\usepackage{colortbl}
\usepackage{makecell}
\usepackage{multirow}
\usepackage{supertabular}

\begin{document}

\newcounter{utterance}

\centering \large Interaction Transcript for game `matchit\_ascii', experiment `different\_grid', episode 2 with claude{-}opus{-}4{-}8{-}azure.
\vspace{24pt}

{ \footnotesize  \setcounter{utterance}{1}
\setlength{\tabcolsep}{0pt}
\begin{supertabular}{c@{$\;$}|p{.15\linewidth}@{}p{.15\linewidth}p{.15\linewidth}p{.15\linewidth}p{.15\linewidth}p{.15\linewidth}}
    \# & $\;$A & \multicolumn{4}{c}{Game Master} & $\;\:$B\\
    \hline

    \theutterance \stepcounter{utterance}  
    & & \multicolumn{4}{p{0.6\linewidth}}{
        \cellcolor[rgb]{0.9,0.9,0.9}{
            \makecell[{{p{\linewidth}}}]{
                \texttt{\tiny{[P1$\langle$GM]}}
                \texttt{Bierzesz udział w grze polegającej na wspólnym zgadywaniu.} \\
\texttt{Celem jest ustalenie, czy ta siatka i inny obrazek, który widzę tylko ja, są takie same.} \\
\texttt{Proszę najpierw opisać swoją siatkę. Następnie ja przedstawię swój opis i będziemy mogli zadawać sobie nawzajem pytania dotyczące siatek, aby ustalić, czy są takie same. Teraz zacznij swój krótki opis siatki od "OPIS:", a następnie podaj opis. Nie dodawaj nic więcej. To jest twoja siatka:} \\
\texttt{X X X X X} \\
\texttt{X ▢ ▢ ▢ ▢} \\
\texttt{X ▢ ▢ ▢ ▢} \\
\texttt{X ▢ ▢ ▢ ▢} \\
\texttt{X ▢ ▢ ▢ ▢} \\
            }
        }
    }
    & \\ \\

    \theutterance \stepcounter{utterance}  
    & \multicolumn{4}{p{0.6\linewidth}}{
        \cellcolor[rgb]{0.9,0.9,0.9}{
            \makecell[{{p{\linewidth}}}]{
                \texttt{\tiny{[P1$\rangle$GM]}}
                \texttt{OPIS: Siatka ma rozmiar 5x5. Cały górny wiersz jest wypełniony znakami X oraz cała lewa kolumna jest wypełniona znakami X. Pozostałe pola (prawy dolny obszar 4x4) są puste. Tworzy to kształt litery L odwróconej — X{-}y wzdłuż górnej krawędzi i lewej krawędzi.} \\
            }
        }
    }
    & & \\ \\

    \theutterance \stepcounter{utterance}  
    & & & \multicolumn{2}{p{0.3\linewidth}}{
        \cellcolor[rgb]{0.9,0.9,0.9}{
            \makecell[{{p{\linewidth}}}]{
                \texttt{\tiny{[GM$|$GM]}}
                \texttt{continue} \\
            }
        }
    }
    & & \\ \\

    \theutterance \stepcounter{utterance}  
    & & \multicolumn{4}{p{0.6\linewidth}}{
        \cellcolor[rgb]{0.9,0.9,0.9}{
            \makecell[{{p{\linewidth}}}]{
                \texttt{\tiny{[GM$\rangle$P2]}}
                \texttt{Bierzesz udział w grze polegającej na wspólnym zgadywaniu.} \\
\texttt{Celem jest ustalenie, czy ta siatka i inny obrazek, który widzę tylko ja, są takie same.} \\
\texttt{Proszę najpierw opisać swoją siatkę. Następnie ja przedstawię swój opis i będziemy mogli zadawać sobie nawzajem pytania dotyczące siatek, aby ustalić, czy są takie same. Teraz zacznij swój krótki opis siatki od "OPIS:", a następnie podaj opis. Nie dodawaj nic więcej. To jest twoja siatka:} \\
\texttt{▢ ▢ X ▢ ▢} \\
\texttt{▢ ▢ X ▢ ▢} \\
\texttt{▢ X X X ▢} \\
\texttt{▢ X X X ▢} \\
\texttt{X X X X X} \\
            }
        }
    }
    & \\ \\

    \theutterance \stepcounter{utterance}  
    & & & \multicolumn{4}{p{0.6\linewidth}}{
        \cellcolor[rgb]{0.9,0.9,0.9}{
            \makecell[{{p{\linewidth}}}]{
                \texttt{\tiny{[GM$\langle$P2]}}
                \texttt{OPIS: Mam siatkę 5x5. Znaki X tworzą kształt przypominający choinkę lub piramidę. W pierwszym wierszu X jest w środkowej kolumnie (3). W drugim wierszu również X w kolumnie 3. W trzecim wierszu X{-}y w kolumnach 2, 3, 4. W czwartym wierszu X{-}y w kolumnach 2, 3, 4. W piątym (ostatnim) wierszu X{-}y we wszystkich pięciu kolumnach. Pozostałe pola są puste.} \\
            }
        }
    }
     \\ \\

    \theutterance \stepcounter{utterance}  
    & & & \multicolumn{2}{p{0.3\linewidth}}{
        \cellcolor[rgb]{0.9,0.9,0.9}{
            \makecell[{{p{\linewidth}}}]{
                \texttt{\tiny{[GM$|$GM]}}
                \texttt{continue} \\
            }
        }
    }
    & & \\ \\

    \theutterance \stepcounter{utterance}  
    & & \multicolumn{4}{p{0.6\linewidth}}{
        \cellcolor[rgb]{0.9,0.9,0.9}{
            \makecell[{{p{\linewidth}}}]{
                \texttt{\tiny{[GM$\rangle$P2]}}
                \texttt{To jest mój OPISOPIS: Siatka ma rozmiar 5x5. Cały górny wiersz jest wypełniony znakami X oraz cała lewa kolumna jest wypełniona znakami X. Pozostałe pola (prawy dolny obszar 4x4) są puste. Tworzy to kształt litery L odwróconej — X{-}y wzdłuż górnej krawędzi i lewej krawędzi.} \\
\texttt{Teraz zadaj pytanie, aby dowiedzieć się o nowych aspektach mojej siatki, które mogą różnić się od twojej siatki. Zacznij od "PYTANIE:" i nie dodawaj nic więcej.} \\
            }
        }
    }
    & \\ \\

    \theutterance \stepcounter{utterance}  
    & & & \multicolumn{4}{p{0.6\linewidth}}{
        \cellcolor[rgb]{0.9,0.9,0.9}{
            \makecell[{{p{\linewidth}}}]{
                \texttt{\tiny{[GM$\langle$P2]}}
                \texttt{PYTANIE: Czy w Twojej siatce ostatni (piąty) wiersz jest wypełniony znakami X we wszystkich pięciu kolumnach, czy tylko w pierwszej kolumnie?} \\
            }
        }
    }
     \\ \\

    \theutterance \stepcounter{utterance}  
    & & & \multicolumn{2}{p{0.3\linewidth}}{
        \cellcolor[rgb]{0.9,0.9,0.9}{
            \makecell[{{p{\linewidth}}}]{
                \texttt{\tiny{[GM$|$GM]}}
                \texttt{continue} \\
            }
        }
    }
    & & \\ \\

    \theutterance \stepcounter{utterance}  
    & & \multicolumn{4}{p{0.6\linewidth}}{
        \cellcolor[rgb]{0.9,0.9,0.9}{
            \makecell[{{p{\linewidth}}}]{
                \texttt{\tiny{[P1$\langle$GM]}}
                \texttt{To jest mój OPISOPIS: Mam siatkę 5x5. Znaki X tworzą kształt przypominający choinkę lub piramidę. W pierwszym wierszu X jest w środkowej kolumnie (3). W drugim wierszu również X w kolumnie 3. W trzecim wierszu X{-}y w kolumnach 2, 3, 4. W czwartym wierszu X{-}y w kolumnach 2, 3, 4. W piątym (ostatnim) wierszu X{-}y we wszystkich pięciu kolumnach. Pozostałe pola są puste.} \\
\texttt{PYTANIE: Czy w Twojej siatce ostatni (piąty) wiersz jest wypełniony znakami X we wszystkich pięciu kolumnach, czy tylko w pierwszej kolumnie?Rozpocznij odpowiedź od "ODPOWIEDŹ:" i nie dodawaj nic więcej.} \\
            }
        }
    }
    & \\ \\

    \theutterance \stepcounter{utterance}  
    & \multicolumn{4}{p{0.6\linewidth}}{
        \cellcolor[rgb]{0.9,0.9,0.9}{
            \makecell[{{p{\linewidth}}}]{
                \texttt{\tiny{[P1$\rangle$GM]}}
                \texttt{ODPOWIEDŹ: W mojej siatce ostatni (piąty) wiersz ma znak X tylko w pierwszej kolumnie, pozostałe pola tego wiersza są puste. Nasze siatki są różne.} \\
            }
        }
    }
    & & \\ \\

    \theutterance \stepcounter{utterance}  
    & & & \multicolumn{2}{p{0.3\linewidth}}{
        \cellcolor[rgb]{0.9,0.9,0.9}{
            \makecell[{{p{\linewidth}}}]{
                \texttt{\tiny{[GM$|$GM]}}
                \texttt{continue} \\
            }
        }
    }
    & & \\ \\

    \theutterance \stepcounter{utterance}  
    & & \multicolumn{4}{p{0.6\linewidth}}{
        \cellcolor[rgb]{0.9,0.9,0.9}{
            \makecell[{{p{\linewidth}}}]{
                \texttt{\tiny{[P1$\langle$GM]}}
                \texttt{Teraz zadaj pytanie, aby dowiedzieć się o nowych aspektach mojej siatki, które mogą różnić się od twojej siatki. Zacznij od "PYTANIE:" i nie dodawaj nic więcej.} \\
            }
        }
    }
    & \\ \\

    \theutterance \stepcounter{utterance}  
    & \multicolumn{4}{p{0.6\linewidth}}{
        \cellcolor[rgb]{0.9,0.9,0.9}{
            \makecell[{{p{\linewidth}}}]{
                \texttt{\tiny{[P1$\rangle$GM]}}
                \texttt{PYTANIE: Czy w Twojej siatce górny wiersz (pierwszy) ma znaki X we wszystkich pięciu kolumnach, czy tylko w środkowej kolumnie (3)?} \\
            }
        }
    }
    & & \\ \\

    \theutterance \stepcounter{utterance}  
    & & & \multicolumn{2}{p{0.3\linewidth}}{
        \cellcolor[rgb]{0.9,0.9,0.9}{
            \makecell[{{p{\linewidth}}}]{
                \texttt{\tiny{[GM$|$GM]}}
                \texttt{continue} \\
            }
        }
    }
    & & \\ \\

    \theutterance \stepcounter{utterance}  
    & & \multicolumn{4}{p{0.6\linewidth}}{
        \cellcolor[rgb]{0.9,0.9,0.9}{
            \makecell[{{p{\linewidth}}}]{
                \texttt{\tiny{[GM$\rangle$P2]}}
                \texttt{ODPOWIEDŹ: W mojej siatce ostatni (piąty) wiersz ma znak X tylko w pierwszej kolumnie, pozostałe pola tego wiersza są puste. Nasze siatki są różne.} \\
\texttt{PYTANIE: Czy w Twojej siatce górny wiersz (pierwszy) ma znaki X we wszystkich pięciu kolumnach, czy tylko w środkowej kolumnie (3)?Rozpocznij odpowiedź od "ODPOWIEDŹ:" i nie dodawaj nic więcej.} \\
            }
        }
    }
    & \\ \\

    \theutterance \stepcounter{utterance}  
    & & & \multicolumn{4}{p{0.6\linewidth}}{
        \cellcolor[rgb]{0.9,0.9,0.9}{
            \makecell[{{p{\linewidth}}}]{
                \texttt{\tiny{[GM$\langle$P2]}}
                \texttt{ODPOWIEDŹ: W mojej siatce górny wiersz ma znak X tylko w środkowej kolumnie (3), pozostałe pola tego wiersza są puste. Nasze siatki są różne.} \\
            }
        }
    }
     \\ \\

    \theutterance \stepcounter{utterance}  
    & & & \multicolumn{2}{p{0.3\linewidth}}{
        \cellcolor[rgb]{0.9,0.9,0.9}{
            \makecell[{{p{\linewidth}}}]{
                \texttt{\tiny{[GM$|$GM]}}
                \texttt{continue} \\
            }
        }
    }
    & & \\ \\

    \theutterance \stepcounter{utterance}  
    & & \multicolumn{4}{p{0.6\linewidth}}{
        \cellcolor[rgb]{0.9,0.9,0.9}{
            \makecell[{{p{\linewidth}}}]{
                \texttt{\tiny{[GM$\rangle$P2]}}
                \texttt{Teraz zadaj pytanie, aby dowiedzieć się o nowych aspektach mojej siatki, które mogą różnić się od twojej siatki. Zacznij od "PYTANIE:" i nie dodawaj nic więcej.} \\
            }
        }
    }
    & \\ \\

    \theutterance \stepcounter{utterance}  
    & & & \multicolumn{4}{p{0.6\linewidth}}{
        \cellcolor[rgb]{0.9,0.9,0.9}{
            \makecell[{{p{\linewidth}}}]{
                \texttt{\tiny{[GM$\langle$P2]}}
                \texttt{PYTANIE: Czy w Twojej siatce lewa kolumna (kolumna 1) jest w całości wypełniona znakami X we wszystkich pięciu wierszach?} \\
            }
        }
    }
     \\ \\

    \theutterance \stepcounter{utterance}  
    & & & \multicolumn{2}{p{0.3\linewidth}}{
        \cellcolor[rgb]{0.9,0.9,0.9}{
            \makecell[{{p{\linewidth}}}]{
                \texttt{\tiny{[GM$|$GM]}}
                \texttt{continue} \\
            }
        }
    }
    & & \\ \\

    \theutterance \stepcounter{utterance}  
    & & \multicolumn{4}{p{0.6\linewidth}}{
        \cellcolor[rgb]{0.9,0.9,0.9}{
            \makecell[{{p{\linewidth}}}]{
                \texttt{\tiny{[P1$\langle$GM]}}
                \texttt{ODPOWIEDŹ: W mojej siatce górny wiersz ma znak X tylko w środkowej kolumnie (3), pozostałe pola tego wiersza są puste. Nasze siatki są różne.} \\
\texttt{PYTANIE: Czy w Twojej siatce lewa kolumna (kolumna 1) jest w całości wypełniona znakami X we wszystkich pięciu wierszach?Rozpocznij odpowiedź od "ODPOWIEDŹ:" i nie dodawaj nic więcej.} \\
            }
        }
    }
    & \\ \\

    \theutterance \stepcounter{utterance}  
    & \multicolumn{4}{p{0.6\linewidth}}{
        \cellcolor[rgb]{0.9,0.9,0.9}{
            \makecell[{{p{\linewidth}}}]{
                \texttt{\tiny{[P1$\rangle$GM]}}
                \texttt{ODPOWIEDŹ: Tak, w mojej siatce lewa kolumna (kolumna 1) jest w całości wypełniona znakami X we wszystkich pięciu wierszach. Nasze siatki są różne.} \\
            }
        }
    }
    & & \\ \\

    \theutterance \stepcounter{utterance}  
    & & & \multicolumn{2}{p{0.3\linewidth}}{
        \cellcolor[rgb]{0.9,0.9,0.9}{
            \makecell[{{p{\linewidth}}}]{
                \texttt{\tiny{[GM$|$GM]}}
                \texttt{continue} \\
            }
        }
    }
    & & \\ \\

    \theutterance \stepcounter{utterance}  
    & & \multicolumn{4}{p{0.6\linewidth}}{
        \cellcolor[rgb]{0.9,0.9,0.9}{
            \makecell[{{p{\linewidth}}}]{
                \texttt{\tiny{[P1$\langle$GM]}}
                \texttt{Teraz zadaj pytanie, aby dowiedzieć się o nowych aspektach mojej siatki, które mogą różnić się od twojej siatki. Zacznij od "PYTANIE:" i nie dodawaj nic więcej.} \\
            }
        }
    }
    & \\ \\

    \theutterance \stepcounter{utterance}  
    & \multicolumn{4}{p{0.6\linewidth}}{
        \cellcolor[rgb]{0.9,0.9,0.9}{
            \makecell[{{p{\linewidth}}}]{
                \texttt{\tiny{[P1$\rangle$GM]}}
                \texttt{PYTANIE: Czy w Twojej siatce trzeci wiersz ma znaki X w kolumnach 2, 3 i 4, przy czym kolumny 1 i 5 są puste?} \\
            }
        }
    }
    & & \\ \\

    \theutterance \stepcounter{utterance}  
    & & & \multicolumn{2}{p{0.3\linewidth}}{
        \cellcolor[rgb]{0.9,0.9,0.9}{
            \makecell[{{p{\linewidth}}}]{
                \texttt{\tiny{[GM$|$GM]}}
                \texttt{continue} \\
            }
        }
    }
    & & \\ \\

    \theutterance \stepcounter{utterance}  
    & & \multicolumn{4}{p{0.6\linewidth}}{
        \cellcolor[rgb]{0.9,0.9,0.9}{
            \makecell[{{p{\linewidth}}}]{
                \texttt{\tiny{[GM$\rangle$P2]}}
                \texttt{ODPOWIEDŹ: Tak, w mojej siatce lewa kolumna (kolumna 1) jest w całości wypełniona znakami X we wszystkich pięciu wierszach. Nasze siatki są różne.} \\
\texttt{PYTANIE: Czy w Twojej siatce trzeci wiersz ma znaki X w kolumnach 2, 3 i 4, przy czym kolumny 1 i 5 są puste?Rozpocznij odpowiedź od "ODPOWIEDŹ:" i nie dodawaj nic więcej.} \\
            }
        }
    }
    & \\ \\

    \theutterance \stepcounter{utterance}  
    & & & \multicolumn{4}{p{0.6\linewidth}}{
        \cellcolor[rgb]{0.9,0.9,0.9}{
            \makecell[{{p{\linewidth}}}]{
                \texttt{\tiny{[GM$\langle$P2]}}
                \texttt{ODPOWIEDŹ: Tak, w mojej siatce trzeci wiersz ma znaki X w kolumnach 2, 3 i 4, a kolumny 1 i 5 są puste. Nasze siatki są różne.} \\
            }
        }
    }
     \\ \\

    \theutterance \stepcounter{utterance}  
    & & & \multicolumn{2}{p{0.3\linewidth}}{
        \cellcolor[rgb]{0.9,0.9,0.9}{
            \makecell[{{p{\linewidth}}}]{
                \texttt{\tiny{[GM$|$GM]}}
                \texttt{continue} \\
            }
        }
    }
    & & \\ \\

    \theutterance \stepcounter{utterance}  
    & & \multicolumn{4}{p{0.6\linewidth}}{
        \cellcolor[rgb]{0.9,0.9,0.9}{
            \makecell[{{p{\linewidth}}}]{
                \texttt{\tiny{[GM$\rangle$P2]}}
                \texttt{Teraz zadaj pytanie, aby dowiedzieć się o nowych aspektach mojej siatki, które mogą różnić się od twojej siatki. Zacznij od "PYTANIE:" i nie dodawaj nic więcej.} \\
            }
        }
    }
    & \\ \\

    \theutterance \stepcounter{utterance}  
    & & & \multicolumn{4}{p{0.6\linewidth}}{
        \cellcolor[rgb]{0.9,0.9,0.9}{
            \makecell[{{p{\linewidth}}}]{
                \texttt{\tiny{[GM$\langle$P2]}}
                \texttt{PYTANIE: Czy w Twojej siatce prawa kolumna (kolumna 5) ma jakiekolwiek znaki X poza górnym wierszem?} \\
            }
        }
    }
     \\ \\

    \theutterance \stepcounter{utterance}  
    & & & \multicolumn{2}{p{0.3\linewidth}}{
        \cellcolor[rgb]{0.9,0.9,0.9}{
            \makecell[{{p{\linewidth}}}]{
                \texttt{\tiny{[GM$|$GM]}}
                \texttt{continue} \\
            }
        }
    }
    & & \\ \\

    \theutterance \stepcounter{utterance}  
    & & \multicolumn{4}{p{0.6\linewidth}}{
        \cellcolor[rgb]{0.9,0.9,0.9}{
            \makecell[{{p{\linewidth}}}]{
                \texttt{\tiny{[P1$\langle$GM]}}
                \texttt{ODPOWIEDŹ: Tak, w mojej siatce trzeci wiersz ma znaki X w kolumnach 2, 3 i 4, a kolumny 1 i 5 są puste. Nasze siatki są różne.} \\
\texttt{PYTANIE: Czy w Twojej siatce prawa kolumna (kolumna 5) ma jakiekolwiek znaki X poza górnym wierszem?Rozpocznij odpowiedź od "ODPOWIEDŹ:" i nie dodawaj nic więcej.} \\
            }
        }
    }
    & \\ \\

    \theutterance \stepcounter{utterance}  
    & \multicolumn{4}{p{0.6\linewidth}}{
        \cellcolor[rgb]{0.9,0.9,0.9}{
            \makecell[{{p{\linewidth}}}]{
                \texttt{\tiny{[P1$\rangle$GM]}}
                \texttt{ODPOWIEDŹ: Nie, w mojej siatce prawa kolumna (kolumna 5) ma znak X tylko w górnym wierszu, a w pozostałych wierszach (2{-}5) jest pusta. Nasze siatki są różne.} \\
            }
        }
    }
    & & \\ \\

    \theutterance \stepcounter{utterance}  
    & & & \multicolumn{2}{p{0.3\linewidth}}{
        \cellcolor[rgb]{0.9,0.9,0.9}{
            \makecell[{{p{\linewidth}}}]{
                \texttt{\tiny{[GM$|$GM]}}
                \texttt{continue} \\
            }
        }
    }
    & & \\ \\

    \theutterance \stepcounter{utterance}  
    & & \multicolumn{4}{p{0.6\linewidth}}{
        \cellcolor[rgb]{0.9,0.9,0.9}{
            \makecell[{{p{\linewidth}}}]{
                \texttt{\tiny{[P1$\langle$GM]}}
                \texttt{Teraz zadaj pytanie, aby dowiedzieć się o nowych aspektach mojej siatki, które mogą różnić się od twojej siatki. Zacznij od "PYTANIE:" i nie dodawaj nic więcej.} \\
            }
        }
    }
    & \\ \\

    \theutterance \stepcounter{utterance}  
    & \multicolumn{4}{p{0.6\linewidth}}{
        \cellcolor[rgb]{0.9,0.9,0.9}{
            \makecell[{{p{\linewidth}}}]{
                \texttt{\tiny{[P1$\rangle$GM]}}
                \texttt{PYTANIE: Czy w Twojej siatce drugi wiersz ma znak X tylko w kolumnie 3, a pozostałe kolumny (1, 2, 4, 5) są puste?} \\
            }
        }
    }
    & & \\ \\

    \theutterance \stepcounter{utterance}  
    & & & \multicolumn{2}{p{0.3\linewidth}}{
        \cellcolor[rgb]{0.9,0.9,0.9}{
            \makecell[{{p{\linewidth}}}]{
                \texttt{\tiny{[GM$|$GM]}}
                \texttt{continue} \\
            }
        }
    }
    & & \\ \\

    \theutterance \stepcounter{utterance}  
    & & \multicolumn{4}{p{0.6\linewidth}}{
        \cellcolor[rgb]{0.9,0.9,0.9}{
            \makecell[{{p{\linewidth}}}]{
                \texttt{\tiny{[GM$\rangle$P2]}}
                \texttt{ODPOWIEDŹ: Nie, w mojej siatce prawa kolumna (kolumna 5) ma znak X tylko w górnym wierszu, a w pozostałych wierszach (2{-}5) jest pusta. Nasze siatki są różne.} \\
\texttt{PYTANIE: Czy w Twojej siatce drugi wiersz ma znak X tylko w kolumnie 3, a pozostałe kolumny (1, 2, 4, 5) są puste?Rozpocznij odpowiedź od "ODPOWIEDŹ:" i nie dodawaj nic więcej.} \\
            }
        }
    }
    & \\ \\

    \theutterance \stepcounter{utterance}  
    & & & \multicolumn{4}{p{0.6\linewidth}}{
        \cellcolor[rgb]{0.9,0.9,0.9}{
            \makecell[{{p{\linewidth}}}]{
                \texttt{\tiny{[GM$\langle$P2]}}
                \texttt{ODPOWIEDŹ: Tak, w mojej siatce drugi wiersz ma znak X tylko w kolumnie 3, a pozostałe kolumny (1, 2, 4, 5) są puste. Nasze siatki są różne.} \\
            }
        }
    }
     \\ \\

    \theutterance \stepcounter{utterance}  
    & & & \multicolumn{2}{p{0.3\linewidth}}{
        \cellcolor[rgb]{0.9,0.9,0.9}{
            \makecell[{{p{\linewidth}}}]{
                \texttt{\tiny{[GM$|$GM]}}
                \texttt{continue} \\
            }
        }
    }
    & & \\ \\

    \theutterance \stepcounter{utterance}  
    & & \multicolumn{4}{p{0.6\linewidth}}{
        \cellcolor[rgb]{0.9,0.9,0.9}{
            \makecell[{{p{\linewidth}}}]{
                \texttt{\tiny{[GM$\rangle$P2]}}
                \texttt{Teraz podejmij decyzję. Jak myślisz: czy Twoja siatka i druga opisana siatka są takie same? Napisz "DECYZJA: ta sama siatka" jeśli uważasz, że to ta sama siatka, albo "DECYZJA: różne siatki" jeśli uważasz, że to różne siatki. Nie dodawaj niczego więcej.} \\
            }
        }
    }
    & \\ \\

    \theutterance \stepcounter{utterance}  
    & & & \multicolumn{4}{p{0.6\linewidth}}{
        \cellcolor[rgb]{0.9,0.9,0.9}{
            \makecell[{{p{\linewidth}}}]{
                \texttt{\tiny{[GM$\langle$P2]}}
                \texttt{DECYZJA: różne siatki} \\
            }
        }
    }
     \\ \\

    \theutterance \stepcounter{utterance}  
    & & & \multicolumn{2}{p{0.3\linewidth}}{
        \cellcolor[rgb]{0.9,0.9,0.9}{
            \makecell[{{p{\linewidth}}}]{
                \texttt{\tiny{[GM$|$GM]}}
                \texttt{continue} \\
            }
        }
    }
    & & \\ \\

    \theutterance \stepcounter{utterance}  
    & & & \multicolumn{2}{p{0.3\linewidth}}{
        \cellcolor[rgb]{0.9,0.9,0.9}{
            \makecell[{{p{\linewidth}}}]{
                \texttt{\tiny{[GM$|$GM]}}
                \texttt{success} \\
            }
        }
    }
    & & \\ \\

    \theutterance \stepcounter{utterance}  
    & & \multicolumn{4}{p{0.6\linewidth}}{
        \cellcolor[rgb]{0.9,0.9,0.9}{
            \makecell[{{p{\linewidth}}}]{
                \texttt{\tiny{[P1$\langle$GM]}}
                \texttt{ODPOWIEDŹ: Tak, w mojej siatce drugi wiersz ma znak X tylko w kolumnie 3, a pozostałe kolumny (1, 2, 4, 5) są puste. Nasze siatki są różne.} \\
\texttt{Teraz podejmij decyzję. Jak myślisz: czy Twoja siatka i druga opisana siatka są takie same? Napisz "DECYZJA: ta sama siatka" jeśli uważasz, że to ta sama siatka, albo "DECYZJA: różne siatki" jeśli uważasz, że to różne siatki. Nie dodawaj niczego więcej.} \\
            }
        }
    }
    & \\ \\

    \theutterance \stepcounter{utterance}  
    & \multicolumn{4}{p{0.6\linewidth}}{
        \cellcolor[rgb]{0.9,0.9,0.9}{
            \makecell[{{p{\linewidth}}}]{
                \texttt{\tiny{[P1$\rangle$GM]}}
                \texttt{DECYZJA: różne siatki} \\
            }
        }
    }
    & & \\ \\

    \theutterance \stepcounter{utterance}  
    & & & \multicolumn{2}{p{0.3\linewidth}}{
        \cellcolor[rgb]{0.9,0.9,0.9}{
            \makecell[{{p{\linewidth}}}]{
                \texttt{\tiny{[GM$|$GM]}}
                \texttt{continue} \\
            }
        }
    }
    & & \\ \\

    \theutterance \stepcounter{utterance}  
    & & & \multicolumn{2}{p{0.3\linewidth}}{
        \cellcolor[rgb]{0.9,0.9,0.9}{
            \makecell[{{p{\linewidth}}}]{
                \texttt{\tiny{[GM$|$GM]}}
                \texttt{success} \\
            }
        }
    }
    & & \\ \\

\end{supertabular}
}

\end{document}
